\section{Objetivos}

O objetivo principal deste estudo é desenvolver um modelo preditivo utilizando técnicas de análise de sobrevivência para prever o comportamento de compra dos clientes em um ambiente de e-commerce. O foco é modelar o tempo até a próxima compra ou o abandono do cliente, a fim de fornecer insights valiosos para otimização de estratégias de marketing e aumento da retenção de clientes.

Os objetivos específicos para o desenvolvimento do trabalho são:
\begin{itemize}
    \item Obter e preparar os dados do conjunto ``E-commerce Customer Behavior and Purchase Datase'';
    \item Realizar o pré-processamento dos dados, incluindo a limpeza, transformação e seleção das variáveis relevantes para a análise;
    \item Aplicar a análise de sobrevivência, utilizando o modelo de Cox e os algoritmos árvore de decisão e floresta aleatória, para identificar os fatores que influenciam o comportamento de compra e o tempo até a próxima compra ou abandono;
    \item Avaliar a performance do modelo gerado, utilizando métricas adequadas para análise de sobrevivência, como o valor p, a curva de Kaplan-Meier e o risco relativo.
    \item Interpretar os resultados para oferecer recomendações práticas para as empresas de e-commerce, visando melhorar as estratégias de marketing, aumentar a fidelização e reduzir a rotatividade de clientes.
\end{itemize}

